\section{Critical Reflection}
\label{sec:critical_reflection}

In order to appropriately contextualise the empirical findings, this section provides a critical analysis of the technical trade-offs and methodological limitations encountered during the development of the \texttt{DetectorNX} project. 
While the results validate the efficiency of the spiking paradigm, several architectural and environmental constraints persist that define the frontier of current \gls{snn} research.

\subsection{Hardware-Paradigm Inequity: The \gls{lif} Simulation Penalty}
The most significant limitation of this evaluation is the hardware-mismatch penalty inherent in benchmarking an event-driven \gls{snn} on synchronous \gls{gpu} hardware. 
Utilizing the SpikingJelly library's multi-step execution paradigm (\texttt{step\_mode='m'}) \cite{fang2023spikingjelly}, the $T=10$ temporal window is processed as a series of fused-kernel iterations. The \gls{gpu} must explicitly update the membrane potential and reset state for every neuron in every layer, regardless of whether a spike is generated. 

This ``simulated time'' approach nullifies the asynchronous, event-driven nature of the \gls{snn}, resulting in the $5.4\times$ throughput penalty documented in Section \ref{subsec:exp_latency}. 
To achieve a truly definitive comparison, the architecture should be deployed on non-von Neumann silicon such as the IBM TrueNorth \cite{merolla2014million} or Intel Loihi \cite{davies2018loihi};
only on such hardware can the $94.33\%$ power savings be realized in a live, asynchronous context.

\subsection{The Temporal Quantization Ceiling}
The adoption of a fixed temporal window of $T=10$ timesteps established a global precision ceiling that was particularly apparent in the quantization error analysis (Section \ref{subsec:quantization_error}). 
As established in early spiking vision literature \cite{THORPE2001715}, rate coding over discrete steps creates a $1/T$ ($10\%$) quantization error in the feature maps. 

This representational limit prevents the \gls{snn} from achieving the infinite sub-pixel coordinate refinement of the continuous \gls{cnn} baseline. To break this fixed-precision trade-off, future trials should explore the following approaches:
\begin{itemize}
    \item \textbf{Adaptive Timestepping:} A technique where the model dynamically increases temporal resolution for complex objects.
    \item \textbf{\gls{ttfs}:} An encoding method that utilizes exact spike timing; information is represented by the arrival time of the first spike rather than the firing frequency \cite{ttfs_encoding}.
    \item \textbf{\gls{roc}:} An approach that represents intensities via relative firing orders and sub-millisecond arrival times rather than discrete spike counts \cite{THORPE2001715}.
\end{itemize}   

\subsection{The Detection Head Mismatch}
To enforce strict architectural parity and variable isolation, this project utilized a standard continuous-valued detection head for both models. 
While this ensured a clean empirical comparison of the backbones, it created an output stability gap. The observed ``Hunting Behaviour'' (Section \ref{subsec:temporal_evidence}) is a positive indicator of the \gls{snn}'s high-frequency temporal sensitivity, proving the backbone is successfully reflecting the sub-millisecond jitter of the raw event stream.

However, the static, float-based head was too slow to resolve this highly reactive evidence, leading to the unstable bounding boxes observed in complex scenes. 
The use of a shared head was a deliberate methodological constraint to prevent the introduction of a ``black box'' variable, but it highlights the need for an \gls{snn}-native detection head, such as a membrane-regression head, to stabilize the model's sensitive temporal probing into a smoothed output suitable for \gls{adas} and other applications in this domain.

\subsection{Dataset Domain Gap and Strategic Selection}
Finally, the choice of the \gls{etram} dataset \parencite{Verma_2024_CVPR} introduces a significant domain gap. 
Because the frames are captured from a static roadside perspective ($\sim$6\,m height), they lack the \textbf{ego-motion artifacts}, such as motion blur and rapid optic flow, that define a vehicle-mounted sensor's environment. In real-world autonomous deployments, cameras are mounted directly to moving vehicles; therefore, it is imperative that a perception system is trained to process and filter such dynamic artifacts.

However, the selection of \gls{etram} over alternatives like Prophesee GEN1 was a deliberate methodological trade-off; \gls{etram} was prioritized for its manual label fidelity (human-assigned labels) and \gls{hd} resolution, providing a gold standard benchmark that was essential for the architectural validation of the \texttt{ULTRA} backbone. 
While the static mounting is a weakness, the quality of the annotations ensured that the $94.9\%$ precision parity was a reflection of the architecture's capacity, rather than an artifact of noisy, automated labels.
Furthermore, the \gls{etram} study also provided a custom pre-processing pipeline optimized towards the specific dataset \cite{Verma_2024_CVPR}, which served as the foundation for the \texttt{REPLICA} pipeline detailed in Section \ref{sec:data_preprocessing}. Leveraging this existing infrastructure significantly accelerated development, ensuring the project's ambitious scope could be realized within the constrained timeframe.
